\bookchapter{บทที่ 4: ติดตั้งและจัดการแอป — “ชุดทักษะ” ของโทรศัพท์}

\chapterintro{App ทำให้โทรศัพท์ทำอะไรก็ได้ แต่เราก็ต้องจัดการมันให้ดี เหมือนที่คุณจัดโต๊ะทำงานและมอบสิทธิ์ให้พนักงานใหม่ การติดตั้งแอปก็ควรมีกระบวนการมาตรฐาน}

\booksection{สโตร์ทางการ: “ซูเปอร์มาร์เก็ตแอป” ที่ปลอดภัยที่สุด}

\begin{itemize}
  \item \textbf{แอนดรอยด์}: Google Play (ต่างประเทศ) หรือสโตร์ในตัวของแต่ละแบรนด์ (เช่น AppGallery ของ Huawei, สโตร์ของ Xiaomi)
  \item \textbf{iOS}: มีแค่ App Store
  \item \textbf{HarmonyOS}: AppGallery ของ Huawei
\end{itemize}

ความหมายด้านความปลอดภัยหลักๆ: ผ่านการตรวจสอบ การอัปเดตมีการรับประกัน การชำระเงินปลอดภัย แนวทางปฏิบัติ:\textbf{ดาวน์โหลดจากสโตร์ทางการเท่านั้น!}

\booksection{โหมดบริสุทธิ์/โหมดป้องกัน: เปิด “กำแพงกันไวรัส” ของโทรศัพท์}

แบรนด์จีนกระแสหลัก (Huawei, Xiaomi, OPPO, vivo, Honor) ติดตั้งมาให้ทั้งหมด\textbf{โหมดบริสุทธิ์}หรือ\textbf{โหมดป้องกัน}ซึ่งเป็นสวิตช์ความปลอดภัยระดับระบบ\textbf{แนะนำอย่างยิ่งให้ผู้ใช้ทุกคนเปิดไว้}.

\booksubsection{มันคืออะไรกันแน่?}

โหมดบริสุทธิ์ (หรือเรียกว่า “โหมดป้องกัน”)\textbf{มีฟังก์ชันหลักคือขวางการติดตั้ง App จากภายนอก}: เมื่อเปิดใช้งาน ระบบจะ\textbf{อนุญาตให้ติดตั้ง App จากสโตร์ทางการเท่านั้น}ไฟล์ติดตั้งที่ดาวน์โหลดผ่านเบราว์เซอร์ คิวอาร์โค้ด แอปแชท ฯลฯ จะถูกบล็อกทันที ตัดปัญหาที่ต้นทาง

\booksubsection{เหมาะกับใคร}

\begin{itemize}
  \item \textbf{ผู้สูงอายุ}: แยก "โฆษณาทางการ" กับ "ของปลอม" ไม่ออก โหมดบริสุทธิ์คือเครื่องรางที่มั่นคงที่สุด
  \item \textbf{ผู้ใช้ที่ไม่แคร์เรื่องสิทธิ์ของแอป}: ไม่ได้ต้องติดตั้งแอปจำเพาะ เวอร์ชันทดสอบ หรือเกมเซิร์ฟเวอร์ต่างประเทศ
  \item \textbf{ทุกคนที่อยากเพิ่มความปลอดภัยให้อุปกรณ์}: แม้แต่สายไฟแอนด์ เปิดไว้ก็ยังดีกว่าปิด — พอมีความจำเป็นชั่วคราวค่อยปิดทีหลังก็ได้
\end{itemize}

\booksubsection{เส้นทางเปิดใช้งาน (ยกตัวอย่าง Huawei/Xiaomi)}

\begin{itemize}
  \item \textbf{Huawei/HarmonyOS}: [ตั้งค่า] > [ระบบและการอัปเดต] >\textbf{[โหมดบริสุทธิ์]}แล้วเลือก "เปิด"
  \item \textbf{Xiaomi}: [ตั้งค่า] > [การคุ้มครองความเป็นส่วนตัว] >\textbf{[คุ้มครองความเป็นส่วนตัว]}หรือ [ตั้งค่า] > [ตั้งค่าแอป] > [การติดตั้งแอป] ปิด "ติดตั้งแอปจากแหล่งที่ไม่รู้จัก"
  \item \textbf{OPPO/vivo}: [ตั้งค่า] > [ความปลอดภัย] >\textbf{[ติดตั้งแอปจากแหล่งภายนอก]}แล้วปิดสิทธิ์ทีละแอป
\end{itemize}

\booksubsection{กลไกตรวจสอบความปลอดภัยของสโตร์}

สโตร์แอปไม่ใช่ “ของประดับ” ข้างหลังมีสายการผลิตด้านความปลอดภัยครบวงจร:

\begin{itemize}
  \item \textbf{สแกนไวรัส}: ก่อนแอปจะขึ้นชั้นวาง จะถูกสแกนร่วมกันจากเอนจินกำจัดไวรัสหลายตัว เพื่อจำแนกม้าโทรจันและลักษณะไวรัสที่รู้จักกันแล้ว
  \item \textbf{ตรวจสอบสิทธิ์}: สโตร์จะตรวจว่าสิทธิ์ที่แอปขอ "สมเหตุสมผล" ไหม — แอปไฟฉายจะขอสิทธิ์สมุดโทรศัพท์ทำไม? ปฏิเสธหรือถอดออกจากชั้นทันที
  \item \textbf{ยืนยันตัวนักพัฒนา}: การขึ้นชั้นวางต้องใช้ตัวตนองค์กร บัตรประชาชนบุคคล หรือบัตรเครดิตยืนยัน เกิดปัญหาก็ตามหาตัวคนได้ ไม่ใช่ "แอปนิรนาม"
  \item \textbf{เฝ้าติดตามต่อเนื่อง}: แอปที่ขึ้นชั้นวางแล้วก็ยังถูกสแกนซ้ำเรื่อยๆ พฤติกรรมประสงค์ร้ายที่พบใหม่จะถูกถอดออกจากชั้นวางทันที
\end{itemize}

\booksection{แหล่งภายนอกกับ “ปรสิตอิเล็กทรอนิกส์”: เสี่ยงสูงมาก!}

ดาวน์โหลดแอปจากเว็บหรือลิงก์แปลกๆ อาจแฝงไวรัส ขโมยความเป็นส่วนตัว แอบรูดเงิน ที่นี่เราขอเสนอแนวคิดใหม่เป็นพิเศษ:\textbf{“ปรสิตอิเล็กทรอนิกส์”}.

\textbf{“ปรสิตอิเล็กทรอนิกส์” คืออะไร?}\par

หมายถึงพวก\textbf{ซอฟต์แวร์ “เวอร์ชันแตก” ที่ถูกฝังม้าโทรจัน}— คุณคิดว่า “ประหยัดเงิน” แต่จริงๆ คือ “โดนกับดัก” รูปแบบที่พบบ่อย:

\begin{itemize}
  \item ซอฟต์แวร์เสียเงินแบบ "เวอร์ชันแตก" "เวอร์ชันฟรี" "ตัดโฆษณา"
  \item เครื่องมือสีเทาอย่าง "โกงไลก์" "ชิงอั่งเปา" "โปรแกรมก่อกวนเกม"
  \item แอป "เวอร์ชันภายใน" ที่ "โฆษณาเล็กๆ" ในกลุ่มแชทแนะนำ
\end{itemize}

โปรแกรมพวกนี้เผื่อใช้ได้ แต่ข้างหลังอาจกำลัง\textbf{ขโมยสมุดโทรศัพท์ อ่าน SMS ดักฟังไมโครโฟน ขุดคริปโตในพื้นหลัง}โหมดบริสุทธิ์คือ\textbf{แนวป้องกันด่านแรก}ในการรับมือพวกมัน — ไม่ให้คุณติดตั้งตั้งแต่แรก

\booksubsection{ข้อเสนอแนะด่านสอง:}

\begin{itemize}
  \item \textbf{สแกนไวรัสทั้งเครื่องเป็นประจำ}: ใช้ศูนย์ความปลอดภัยในตัวโทรศัพท์ (เช่น Phone Manager ของ Huawei, Security ของ Xiaomi) สแกนเดือนละครั้ง
  \item \textbf{ระมัดระวังการมอบสิทธิ์ที่อ่อนไหว}: สมุดโทรศัพท์ ข้อความ SMS ตำแหน่ง กล้อง ไมโครโฟน — ห้า "สิทธิ์เสี่ยงสูง" นี้ต้องตรวจพิจารณาอย่างเข้มงวด
  \item \textbf{เฝ้าดูการใช้แบตและดาต้าที่ผิดปกติ}: โทรศัพท์ร้อนฉับพลัน แบตหมดไว ดาต้าวิ่งกลางดึก — แปดทบสิบว่าพื้นหลังแอบเลี้ยง "ปรสิต" ไว้
\end{itemize}

ถ้าหาทางเลี่ยงไม่ได้จริงๆ แล้วต้องติดตั้งแอปภายนอก:

\begin{enumerate}
  \item มั่นใจว่าแหล่งที่มาน่าเชื่อถืออย่างเด็ดขาด
  \item ก่อนติดตั้งอ่านสิทธิ์ที่ขอให้ละเอียด
  \item หลังติดตั้งใช้แอปความปลอดภัยสแกนตาม
\end{enumerate}

คำแนะนำสุดท้าย: ในการตั้งค่าระบบ ปิดสิทธิ์ “ติดตั้งแอปจากแหล่งที่ไม่รู้จัก” ซะ

\booksection{ติดตั้ง อัปเดต ถอนการติดตั้ง}

\begin{itemize}
  \item \textbf{การติดตั้ง}: กด "รับ" หรือ "ติดตั้ง" ในสโตร์ พอติดตั้งเสร็จเปิดครั้งแรกมันจะถามสิทธิ์ที่ต้องการ
  \item \textbf{การอัปเดต}:\textbf{ต้องอัปเดต!}ส่วนใหญ่เพื่อแก้ช่องโหว่ด้านความปลอดภัย แนะนำเปิดอัปเดตอัตโนมัติ
  \item \textbf{การถอนการติดตั้ง}: กดค้างที่ไอคอนลากไปโซนถอนการติดตั้ง ถ้าอยากเก็บกวาดให้สะอาดขึ้น ไปที่ [ตั้งค่า] > [จัดการแอป] ค้นหาแล้วถอนการติดตั้ง
\end{itemize}

\booksection{การจัดการสิทธิ์: คุณคือ “เจ้านาย” ของแอป}

แอปขอสิทธิ์ (เช่นสมุดโทรศัพท์ ตำแหน่ง) บางทีก็สมเหตุสมผล (แผนที่ต้องใช้ตำแหน่ง) บางทีก็ไม่สมเหตุสมผล (แอปไฟฉายจะขอสมุดโทรศัพท์ทำไม)

\begin{itemize}
  \item \textbf{ตัวเลือกสิทธิ์}: "อนุญาตตลอด" "อนุญาตเฉพาะขณะใช้งาน" "ห้าม" "ถามทุกครั้ง"
  \item \textbf{หลักการจัดการ}:
  \begin{enumerate}
    \item ลดให้น้อยที่สุด: ไม่จำเป็นก็ไม่ให้
    \item ตรวจสอบหลังใช้งาน: เข้าไปดูในการตั้งค่าสิทธิ์เป็นระยะว่าใคร "ใช้เกินขอบเขต"
    \item ปรับตามบริบท: เหมือนแอปเรียกรถ ให้สิทธิ์ตำแหน่งเฉพาะตอนใช้งาน
  \end{enumerate}
\end{itemize}

\begin{tcbwarn}{สิทธิ์ที่ต้องเฝ้าระวังเป็นพิเศษ}
\begin{itemize}
  \item สมุดโทรศัพท์/ข้อความ SMS: อาจถูกใช้ส่งสแปม
  \item ตำแหน่งเรียลไทม์: อาจถูกใช้ตามรอยเส้นทางการเดินทาง
  \item กล้อง/ไมโครโฟน: อาจถูกใช้ถ่ายภาพและอัดเสียงแอบแอง
  \item การเข้าถึงไฟล์: อาจอ่านไฟล์ส่วนตัวของคุณ
\end{itemize}

\end{tcbwarn}

เคยสงสัยไหมว่าไฟล์ติดตั้งที่ดาวน์โหลดจากสโตร์นั่น ข้างในมีอะไรกันแน่? การเข้าใจโครงสร้างของไฟล์ติดตั้งช่วยให้คุณเข้าใจว่าแอปทำงานอย่างไร และรู้ว่าต้องไล่ตรวจจากไหนเวลาเจอปัญหา

\booksubsection{โครงสร้างไฟล์ APK ของแอนดรอยด์}

APK (Android Package Kit) แท้จริงคือไฟล์บีบอัด ZIP แตกไฟล์ออกจะเห็นโฟลเดอร์และไฟล์แกนหลักดังนี้:

\begin{itemize}
  \item \textbf{AndroidManifest.xml}: "บัตรประชาชน" และ "คู่มือ" ของแอป บันทึกชื่อแพ็กเกจ เลขเวอร์ชัน สิทธิ์ที่ต้องการ เวอร์ชันระบบที่รองรับ Activity ทางเข้า ฯลฯ ระบบจะอ่านไฟล์นี้ก่อนอย่างอื่นเวลาติดตั้งแอป
  \item \textbf{classes.dex}: โค้ดแกนหลักของแอป ซอร์สโค้ด Java/Kotlin คอมไพล์แล้วกลายเป็นไบต์โค้ด Dalvik เก็บไว้ในไฟล์นี้ แอปหนึ่งตัวอาจมีหลายไฟล์ dex (เช่น classes2.dex, classes3.dex)
  \item \textbf{resources.arsc}: ตารางดัชนีทรัพยากรหลังคอมไพล์ บันทึก ID และความสัมพันธ์ของทรัพยากรทั้งหมด
  \item \textbf{โฟลเดอร์ res/}: โฟลเดอร์ทรัพยากร รวมรูปภาพ ไฟล์เลย์เอาต์ สตริงข้อความ ค่าสี ฯลฯ รูปคนละความละเอียดเก็บในโฟลเดอร์ drawable คนละโฟลเดอร์ (เช่น drawable-hdpi, drawable-xhdpi)
  \item \textbf{โฟลเดอร์ assets/}: โฟลเดอร์ทรัพยากรดิบ เก็บทรัพยากรที่ไม่ผ่านการคอมไพล์ เช่นไฟล์เว็บ ฟอนต์ เสียง ฯลฯ แอปอ่านผ่าน AssetManager ขณะทำงานได้
  \item \textbf{โฟลเดอร์ lib/}: โฟลเดอร์ไลบรารีในเครื่อง เก็บไฟล์ .so ที่คอมไพล์แล้ว (ผลจากโค้ด C/C++) แยกตามสถาปัตยกรรม CPU (เช่น arm64-v8a, armeabi-v7a) เกมและแอปที่ต้องคำนวณหนักมักมีโฟลเดอร์นี้
  \item \textbf{โฟลเดอร์ META-INF/}: โฟลเดอร์ข้อมูลลายเซ็น มี MANIFEST.MF, CERT.SF, CERT.RSA สามไฟล์ ไว้ตรวจสอบความครบถ้วนและลายเซ็นของ APK
\end{itemize}

\booksubsection{โครงสร้างไฟล์ IPA ของ iOS}

IPA (iOS App Store Package) เช่นกันคือไฟล์บีบอัด โครงสร้างต่างกันเล็กน้อย:

\begin{itemize}
  \item \textbf{โฟลเดอร์ Payload/}: ที่อยู่ของเนื้อหาแกนหลัก ข้างในมีโฟลเดอร์ .app หนึ่งอัน รวมทุกอย่างของแอป
  \item \textbf{Info.plist}: ไฟล์ค่าตั้งของแอป คล้าย AndroidManifest.xml ของแอนดรอยด์ บันทึก ID ของแอป เวอร์ชัน สิทธิ์ที่ต้องการ อุปกรณ์ที่รองรับ ฯลฯ
  \item \textbf{ไฟล์ปฏิบัติการ}: โปรแกรมหลักที่ชื่อเดียวกับโฟลเดอร์ .app เป็นไบนารีฟอร์แมต Mach-O หลังคอมไพล์
  \item \textbf{โฟลเดอร์ Frameworks/}: ไลบรารีไดนามิกและเฟรมเวิร์กที่แอปพึ่งพา
  \item \textbf{ไฟล์ทรัพยากร}: รูปภาพ เสียง วิดีโอ ฯลฯ มักเก็บในโฟลเดอร์ .app
  \item \textbf{โฟลเดอร์ \_CodeSignature/}: ข้อมูลลายเซ็นโค้ด ไว้ตรวจสอบความครบถ้วนและที่มาของแอป
\end{itemize}

\booksubsection{กลไกตรวจสอบลายเซ็น}

ทำไมแอปที่ดาวน์โหลดจากสโตร์ทางการถึงปลอดภัยกว่า? การตรวจสอบลายเซ็นคือหัวใจ

\begin{itemize}
  \item \textbf{ลายเซ็นแบบแอนดรอยด์}: นักพัฒนาใช้กุญแจส่วนตัวเซ็น APK ระบบใช้กุญแจสาธารณะตรวจสอบตอนติดตั้ง ถ้า APK ถูกดัดแปลง (เช่นฝังไวรัส) ลายเซ็นจะใช้ไม่ได้ ระบบจะปฏิเสธการติดตั้ง Google Play ยังบังคับ App Signing โดย Google เก็บรักษากุญแจเซ็น ช่วยคุ้มครองแอปเพิ่มขึ้นอีกชั้น
  \item \textbf{ลายเซ็นแบบ iOS}: ฝั่งแอปเปิลตรวจเข้มกว่า ทุก IPA ต้องมีใบรับรองที่แอปเปิลออกให้เซ็น ไม่งั้นติดตั้งไม่ได้ นี่คือเหตุผลที่ iOS ติดตั้งแอปได้จาก App Store หรือใบรับรององค์กรเท่านั้น แอปที่เซ็นด้วยใบรับรององค์กรมีอายุจำกัด และแอปเปิลเรียกคืนได้ทุกเมื่อ
  \item \textbf{ประโยชน์ของลายเซ็น}:
  \begin{enumerate}
    \item ตรวจสอบที่มา: ยืนยันว่าแอปมาจากนักพัฒนาที่อ้างสิทธิ์จริงๆ
    \item ตรวจสอบความครบถ้วน: มั่นใจว่าแอปไม่ถูกดัดแปลงระหว่างการส่งผ่าน
    \item กันการปลอมตัว: แอปประสงค์ร้ายปลอมเป็นแอปอื่นไม่ได้
  \end{enumerate}
\end{itemize}

\booksubsection{ขนาดไฟล์ติดตั้งกับการปรับแต่ง}

คุณอาจพบว่าขนาดแอปที่สโตร์แสดง กับพื้นที่ที่กินหลังติดตั้งไม่เท่ากัน เพราะว่า:

\begin{itemize}
  \item \textbf{ขนาดตอนดาวน์โหลด}: ขนาด APK/IPA หลังบีบอัด
  \item \textbf{ขนาดตอนติดตั้ง}: ขนาดหลังแตกไฟล์ มักใหญ่กว่าตอนดาวน์โหลดมาก
  \item \textbf{พื้นที่ที่ใช้ขณะทำงาน}: แอปขณะทำงานยังสร้างแคชและไฟล์ข้อมูลเพิ่ม กินพื้นที่มากขึ้นอีก
\end{itemize}

เทคนิคปรับแต่งของแอปยุคใหม่:

\begin{itemize}
  \item \textbf{App Bundle (แอนดรอยด์)}: เทคโนโลยีจัดส่งแบบไดนามิกของ Google Play ตอนดาวน์โหลดผู้ใช้จะได้เฉพาะทรัพยากรที่เหมาะกับอุปกรณ์ตัวเอง (เช่นดาวน์โหลดเฉพาะรูปที่ตรงกับความละเอียดหน้าจอ) ลดขนาดดาวน์โหลดได้มาก
  \item \textbf{App Thinning (iOS)}: เทคนิคคล้ายกัน แอปเปิลจะตัดทรัพยากรที่ไม่จำเป็นออกตามอุปกรณ์อัตโนมัติ
  \item \textbf{ดาวน์โหลดตามต้องการ}: เกมและแอปใหญ่ๆ สามารถดาวน์โหลดแพ็กเกจแกนก่อน แล้วค่อยดาวน์โหลดแพ็กเกจทรัพยากรทีหลัง
\end{itemize}

\begin{tcbinfo}{เกร็ดความรู้: ทำไมแอปบางตัวไฟล์ติดตั้งใหญ่มาก?}
แอปเกมมักใหญ่ที่สุด เพราะมีภาพติดความละเอียดสูง โมเดล 3D และไฟล์เสียงจำนวนมาก แอปโซเชียลก็ค่อนข้างใหญ่ เพราะต้องรองรับโมดูลฟังก์ชันหลากหลาย ส่วนแอปเครื่องมือมักเล็กที่สุด ถ้าคุณพบแอปเครื่องมือตัวไหนใหญ่ผิดปกติ ต้องระวังว่าแอบแถมฟังก์ชันที่ไม่จำเป็นมาหรือเปล่า

\end{tcbinfo}

\begin{tcbpractice}{ฝึกปฏิบัติ: ตรวจสอบสิทธิ์ของ 3 แอป}
เลือกแอปที่ใช้บ่อยสัก 3 ตัว (เช่นโซเชียล ช้อปปิ้ง เครื่องมือ) เข้าไปที่หน้าจัดการสิทธิ์แอปใน [ตั้งค่า] ไล่ตรวจทีละรายการ ปิดสิทธิ์ที่อ่อนไหวที่ไม่จำเป็น หรือเปลี่ยนเป็น “อนุญาตเฉพาะขณะใช้งาน”

\end{tcbpractice}
